\begin{abcliste}
\abc
In der Ruhelage ist die Federkraft gleich der Gewichtskraft, $k\,x_0 = m\,g$:
\begin{align}
 k &= \kFF = \frac{\mK\cdot\g}{\xo} = \kF \approx \result{\kFP{0}{2}}
\end{align}

\abc
Nullniveau der Lageenergie: der tiefste Punkt, $s$ unter dem Start. Am Start ruht der Klotz und die Feder ist entspannt (nur Lageenergie), im tiefsten Punkt ruht er wieder (nur Spannenergie):
\begin{align}
 m\,g\,s &= \frac{1}{2}\,k\,s^2 \\
 s &= \frac{2\,m\,g}{k} = 2\,x_0
\end{align}
\begin{align}
 s &= \sTF = 2\cdot\xo = \sT \approx \result{\sTP{-2}{2}}
\end{align}
Der Klotz sinkt doppelt so weit wie bis zur Ruhelage.

\abc
Der Klotz wird schneller, solange die Gewichtskraft grösser ist als die Federkraft, also bis zur Ruhelage $x_0$ unter dem Start. Dort ist er am schnellsten. Energiebilanz vom Start bis dorthin (Nullniveau jetzt in der Ruhelage):
\begin{align}
 m\,g\,x_0 &= \frac{1}{2}\,k\,x_0^2 + \frac{1}{2}\,m\,v^2
\end{align}
Mit $k\,x_0 = m\,g$ ist $\frac{1}{2}\,k\,x_0^2 = \frac{1}{2}\,m\,g\,x_0$, also $\frac{1}{2}\,m\,v^2 = \frac{1}{2}\,m\,g\,x_0$:
\begin{align}
 v_\mathrm{max} &= \vmF = \sqrt{\g\cdot\xo} = \vm \approx \result{\vmP{0}{2}}
\end{align}

\abc
Im tiefsten Punkt ist der Klotz einen Augenblick in Ruhe: \result{v = 0}, ein Umkehrpunkt. Die ganze Energie steckt in der Feder. Die Feder ist dort um $2\,x_0$ gedehnt und zieht doppelt so stark wie in der Ruhelage:
\begin{align}
 F_\mathrm{F} &= \FTF = \FT \approx \result{\FTP{0}{2}} = 2\cdot F_\mathrm{G}, \qquad F_\mathrm{G} = \FGF = \FG
\end{align}
Die resultierende Kraft zeigt nach oben und ist so gross wie die Gewichtskraft: Der Klotz bleibt nicht stehen, sondern wird wieder nach oben beschleunigt. Er schwingt zwischen dem Start und dem tiefsten Punkt hin und her, um die Ruhelage $x_0$ herum. Die Ruhelage ist nicht der tiefste Punkt, sondern der Ort der grössten Geschwindigkeit.
\end{abcliste}
